\documentclass[default]{aastex701}
\usepackage{amsmath}

% | 模式           | 用途   | 特点          |
% | ------------ | ---- | ----------- |
% | `manuscript` | 审稿   | 单栏、双倍行距、最易读 |
% | `preprint`   | 内部流转 | 单栏、较大字号     |
% | `default`    | 默认   | 单栏、10pt     |
% | `twocolumn`  | 接近出版 | 双栏、最终版式     |
% | 选项               | 作用        |
% | ---------------- | --------- |
% | `linenumbers`    | 显示行号      |
% | `trackchanges`   | 标记修改（增删改） |
% | `anonymous`      | 隐藏作者信息    |
% | `longauthorlist` | 强制完整作者列表  |

\newcommand{\vdag}{(v)^\dagger}
\newcommand\aastex{AAS\TeX}
\newcommand\latex{La\TeX}
\begin{document}

\title{Star Catalog Position and Proper Motion Corrections in Natural Satellite Astrometry}

\author{Zhe Yang}
\affiliation{Shanghai University Of Engineering Science}
\email{zhejones@gmail.com}

% \author{H. Y. Zhang}
% \affiliation{Shanghai University Of Engineering Science}
% \email{123456@gmail.com}

% \author{Kai Tang}
% \affiliation{Shanghai Astronomical Observatory}
% \email{123456@gmail.com}

% \correspondingauthor{H Y zhang}
% \email{you9r.email@institution.edu}

% \collaboration{all}{The Terra Mater collaboration}

\begin{abstract}

Accurate orbit determination of minor bodies relies critically on the quality and long-term consistency of optical astrometry. Over more than a century, asteroid observations have been reduced with a wide variety of stellar reference catalogs whose systematic errors propagate directly into orbital solutions. Previous debiasing efforts, most notably those based on 2MASS and Gaia DR2, substantially reduced regional systematics in historical star catalogs. The third Gaia data release (Gaia DR3), however, provides significantly improved astrometric precision and a longer time baseline, enabling a more stable and homogeneous reference frame for re-evaluating historical catalog biases.

In this work, we construct a unified correction framework for stellar position and proper motion systematics using Gaia DR3 as the reference catalog. Seventeen historical and modern astrometric catalogs spanning photographic, meridian-circle, CCD, and space-based surveys are re-analyzed within a consistent ICRS/J2000 realization. Systematic offsets in position and proper motion are estimated on HEALPix tiles and expressed as time-dependent correction vectors applicable to optical minor-body observations. Particular attention is given to reference-frame transformations, epoch alignment, and cross-matching stability, which are formalized into a reproducible algorithmic pipeline.

The resulting correction tables reveal pronounced stratification of systematic structures among different generations of catalogs, including large-amplitude discontinuities inherited from plate-based reductions. A basis-function smoothing scheme is further implemented to mitigate artificial boundary effects while preserving large-scale systematics. The correction tables and software are publicly released to facilitate consistent reprocessing of historical asteroid astrometry within the Gaia DR3 reference frame.

\end{abstract}


\keywords{\uat{Astrometry}{1648} --- \uat{Astrographic catalogs}{205} --- \uat{Gaia}{1643} --- \uat{satellites}{1089} --- \uat{Natural satellites}{1528}}

\section{Introduction} \label{sec:intro}

Precise orbit determination of Solar System minor bodies depends fundamentally on the long-term consistency of optical astrometric measurements\citep{Milani2010}. Over more than a century, millions of observations have been reduced against a succession of stellar reference catalogs, each constructed under distinct observational strategies, reference frames, and reduction methodologies. Although random measurement uncertainties typically dominate individual observations, regional and catalog-dependent systematic errors introduce correlated biases that can propagate into orbital solutions and long-term ephemeris predictions.

Systematic distortions in star catalogs were first quantified in the context of asteroid astrometry by \citet{Chesley2010}, who demonstrated that regional catalog biases measurably affect orbit fits. Building upon that framework, \citet{Farnocchia2015} incorporated proper-motion information and extended the correction scheme to a broader set of widely used catalogs. With the advent of Gaia DR2, \citet{Eggl2020} re-derived catalog-dependent position and proper-motion corrections using a homogeneous, high-precision reference tied to the ICRS, significantly reducing residual large-scale systematics.

The third Gaia data release (Gaia DR3; \citealt{GaiaDR3}) provides improved astrometric precision, a longer observational baseline, and a more stable global solution compared to Gaia DR2. In particular, the extended time coverage enhances proper-motion accuracy and reduces small-scale reference-frame noise, thereby offering a more robust foundation for re-evaluating historical catalog systematics\citep{Lindegren2021EDR3,Fabricius2021}. Given that debiasing corrections are inherently sensitive to the statistical properties of the adopted reference catalog, a reassessment within the Gaia DR3 realization of the ICRS is warranted.

In this work, we construct a unified and reproducible correction framework for historical and modern astrometric catalogs using Gaia DR3 as the reference. Seventeen catalogs spanning photographic plate surveys, meridian-circle observations, compiled fundamental catalogs, CCD-based surveys, and early space astrometry are re-analyzed in a consistent ICRS/J2000 frame\citep{Ma1998ICRF,Fey2015}. Systematic offsets in stellar positions and proper motions are estimated on equal-area HEALPix tiles and expressed as time-dependent correction vectors applicable to optical minor-body astrometry.

Beyond reproducing prior debiasing strategies, we explicitly formalize reference-frame transformations, epoch alignment procedures, and cross-matching constraints that critically affect stability at sub-arcsecond levels. We further quantify spatial discontinuities inherited from plate-based reductions and introduce a basis-function smoothing scheme to mitigate artificial tile-boundary effects while preserving large-scale systematic structure. The resulting correction tables, released together with the processing pipeline, provide a consistent foundation for reprocessing historical asteroid astrometry within the Gaia DR3 reference frame.


\section{Reference and Astrometric Catalogs} \label{sec:catalogs}

\subsection{Gaia Catalog Series} \label{subsec:gaia}

Since 2013, the Gaia mission has carried out all-sky high-precision astrometric observations, providing positional and photometric measurements for approximately 1.8 billion celestial sources. The second Gaia data release (Gaia DR2) represented the first public release of large-scale five-parameter astrometric solutions, including positions, parallaxes, and proper motions. As a result, the Gaia catalog became a fundamental reference for quantifying systematic errors in external star catalogs and for high-accuracy astrometric studies of Solar System small bodies.

Subsequent releases, Gaia EDR3 and Gaia DR3, adopted the same astrometric parameterization as Gaia DR2 but incorporated substantial improvements in the observational time baseline, spacecraft attitude reconstruction, chromaticity calibration, and instrument modeling. In particular, the astrometric solution of Gaia DR3 is based on approximately 34 months of observations, a significant extension compared to the $\sim$22-month baseline of Gaia DR2. Official validation results \citep{GaiaDR3} indicate that, relative to Gaia DR2, Gaia DR3 achieves a typical reduction of about 20\% in parallax formal uncertainties and more than a factor-of-two improvement in proper motion formal uncertainties, demonstrating a substantial enhancement in overall astrometric precision.

\begin{deluxetable*}{cccc}
\tablecaption{Comparison of Gaia DR1, DR2 and Gaia DR3 astrometric properties 
\label{tab:gaia_comparison}}
\tablehead{
\colhead{Property} & \colhead{Gaia DR1} & \colhead{Gaia DR2} & \colhead{Gaia DR3}}
\startdata
Observation time span       & $\sim$14 months            & $\sim$22 months             & $\sim$34 months                \\
Five-parameter solutions    & $\sim$2 $\times$ 10$^{6}$ (TGAS) & $\sim$1.3 $\times$ 10$^{9}$ & $\sim$5.9 $\times$ 10$^{8}$    \\
Six-parameter solutions     & Not available              & Not available               & $\sim$8.8 $\times$ 10$^{8}$    \\
Proper-motion uncertainty   & $\sim$1 mas yr$^{-1}$      & Reference level             & $\sim$factor 2 improvement     \\
Error distribution          & Quasi-Gaussian (TGAS)      & Quasi-Gaussian              & Quasi-Gaussian (5p $>$ 6p)     \\
Uncertainty underestimation & Present                    & Present                     & Present, better quantified     \\
Systematic noise floor (5p) & $\sim$0.3 mas              & Not quantified              & $\sim$15 $\mu$as               \\
\enddata
\end{deluxetable*}

Using the five-parameter astrometric solutions of Gaia DR2 as a reference, \citet{Eggl2020} constructed a regional star catalog correction model and systematically characterized the dominant structures of systematic biases present in several historical star catalogs as applied to minor planet astrometry. The objective of the present work is to extend the applicability of this framework within the Gaia reference system and to provide a reusable implementation, with an emphasis on consistency checks and engineering-level completion of existing conclusions using a more stable astrometric reference. Given the strong sensitivity of star catalog correction methods to the statistical properties of reference catalog errors, Gaia DR3 is adopted as the reference star catalog in this work.

\subsection{Historical Astrometric Catalogs} \label{subsec:historical}

Optical astrometric observations obtained over different historical periods have relied on a variety of stellar reference catalogs. These catalogs differ in their observation epochs, sky coverage, source density, and astrometric content, reflecting successive stages in the development of ground-based and space-based astrometry. Below, we provide a brief overview of the star catalogs considered in this work (see Table \ref{tab:catalog_overview}).

\begin{deluxetable*}{ccccccccc}
\tablecaption{Summary of astrometric catalogs used in this work
\label{tab:catalog_overview}}
\tablehead{\colhead{Catalog} & \colhead{Magnitude} & \colhead{Star Count} & \colhead{Epoch} & \colhead{Reference Frame} & \colhead{Proper Motion} & \colhead{Coverage} & \colhead{CDS ID} & \colhead{Reference}}
\startdata
AC            & $V \approx 6\text{--}11.0$    & 2,026,347     & plate epoch          & B1950   & NO  & 30\%    & \href{https://cdsarc.cds.unistra.fr/viz-bin/cat/I/154}{I/154}  & \citet{Roeser1990} \\
AGK1          & $V \approx 2.2\text{--}9.1$   & 5,954         & mean epoch           & B1875   & NO  & 3\%     & \href{https://cdsarc.cds.unistra.fr/viz-bin/cat/I/310}{I/310}  & \citet{AGK1} \\
AGK3          & $B,V \approx 6\text{--}11$    & 183,145       & observation          & B1950   & YES & 52\%    & \href{https://cdsarc.cds.unistra.fr/viz-bin/cat/I/61B}{I/61B}  & \citet{AGK3} \\
FK4 catalogue & $V \approx 1\text{--}9.5$     & 1,535         & B1950                & B1950   & YES & 3\%     & \href{https://cdsarc.cds.unistra.fr/viz-bin/cat/I/143}{I/143}  & \citet{FK4} \\
SAO           & $V = 0\text{--}11$            & 258,997       & J2000                & J2000   & YES & 98\%    & \href{https://cdsarc.cds.unistra.fr/viz-bin/cat/I/131A}{I/131A} & \citet{SAO1966} \\
ACRS          & $V \approx 6\text{--}11$      & 320,211       & J2000                & J2000   & YES & 100\%   & \href{https://cdsarc.cds.unistra.fr/viz-bin/cat/I/171}{I/171}  & \citet{ACRS1988} \\
ACT           & $V \approx 6\text{--}11.5$    & 988,758       & J2000                & J2000   & YES & 100\%   & \href{https://cdsarc.cds.unistra.fr/viz-bin/cat/I/246}{I/246}  & \citet{Urban1998ACT} \\
PPM           & $V \approx 1\text{--}11.5$    & 378,910       & J2000                & J2000   & YES & 100\%   & \href{https://cdsarc.cds.unistra.fr/viz-bin/cat/I/146}{I/146}, \href{https://cdsarc.cds.unistra.fr/viz-bin/cat/I/193}{I/193} & \citet{Roeser1991PPM} \\
Yale          & $V \approx -1\text{--}6.5$    & 221,338       & observation          & B1950   & YES & 76\%    & \href{https://cdsarc.cds.unistra.fr/viz-bin/cat/I/141}{I/141}  & \citet{Hoffleit1991BSC} \\
Hipparcos     & $V \approx 1\text{--}12.4$    & 118,218       & J1991.25             & J2000   & YES & 94\%    & \href{https://cdsarc.cds.unistra.fr/viz-bin/cat/I/239}{I/239}  & \citet{ESA1997Hipparcos} \\
Tycho 2       & $V \approx 2\text{--}12.5$    & 2,539,913     & J2000                & J2000   & YES & 100\%   & \href{https://cdsarc.cds.unistra.fr/viz-bin/cat/I/259}{I/259}  & \citet{Hog2000Tycho2} \\  
USNO A2.0     & $B,R = 11\text{--}20$         & 526,280,881   & mean epoch           & J2000   & NO  & 100\%   & \href{https://cdsarc.cds.unistra.fr/viz-bin/cat/I/252}{I/252}  & \citet{Monet1998USNOA2} \\
GSC 1.2       & $V \approx 6\text{--}15$      & 18,825,216    & plate epoch          & J2000   & NO  & 100\%   & \href{https://cdsarc.cds.unistra.fr/viz-bin/cat/I/254}{I/254}  & \citet{Morrison2001GSC12} \\
UCAC 2        & $R \approx 8\text{--}16$      & 48,330,571    & J2000                & J2000   & YES & 88\%    & \href{https://cdsarc.cds.unistra.fr/viz-bin/cat/I/289}{I/289}  & \citet{Zacharias2004UCAC2} \\
UCAC 4        & $R \approx 7\text{--}16$      & 113,780,093   & J2000                & J2000   & YES & 100\%   & \href{https://cdsarc.cds.unistra.fr/viz-bin/cat/I/322A}{I/322A} & \citet{Zacharias2013UCAC4} \\
Gaia DR2      & $G = 3\text{--}21$            & 1,331,909,727 & J2015.5              & J2015.5 & YES & 100\%   & \href{https://cdsarc.cds.unistra.fr/viz-bin/cat/I/345}{I/345}  & \citet{GaiaDR2} \\
Gaia DR1      & $G = 5\text{--}20.7$          & 1,142,685,696 & J2015                & J2015   & NO  & 100\%   & \href{https://cdsarc.cds.unistra.fr/viz-bin/cat/I/337}{I/337}  & \citet{GaiaDR1} \\
\enddata
\tablecomments{Primary literature references correspond to the original catalog publications listed in the CDS entries.}
\end{deluxetable*}
% \begin{itemize}
%     \item \textbf{Astrographic Catalogue:} \label{subsubsec:ac} The Astrographic Catalogue (AC) is one of the earliest large-scale photographic astrometric catalogs, produced within the framework of the international *Carte du Ciel* project initiated at the end of the 19th century. The project employed standardized astrographs to photograph the sky in declination zones assigned to different observatories, with the bulk of the observations carried out between the late 1890s and the first half of the 20th century. The AC comprises several million stellar positions derived from photographic plates, with typical positional accuracies at the few hundred milliarcsecond level. Owing to its very early mean epoch, AC primarily represents the reference framework of early photographic astrometry and provides the longest available astrometric time baseline among the catalogs considered in this work.
%     \item \textbf{Astronomische Gesellschaft Katalog 1:} \label{subsubsec:agk1} AGK1 is a fundamental star catalog based on visual meridian-circle observations, covering the northern sky down to approximately -2 degrees in declination. Its observations were obtained mainly between the late 19th century and the early 20th century. With a limiting magnitude of about 9 and positional accuracies at the arcsecond level, AGK1 represents the reference framework of pre-photographic astrometry. It served as a primary reference catalog for several early photographic surveys, including parts of the Astrographic Catalogue.
%     \item \textbf{Astronomische Gesellschaft Katalog 3:} \label{subsubsec:agk3} AGK3 is a photographic catalog covering the northern sky, developed as the third installment of the AGK series. Its observations were carried out around the late 1950s, and it introduced proper motion information by re-observing earlier AGK stars. With typical positional accuracies of about 0.1 arcseconds and proper motion accuracies at the level of several milliarcseconds per year, AGK3 became a widely adopted reference catalog for ground-based astrometry during the 1970s–1990s. The catalog is strictly tied to the FK4 reference frame.
%     \item \textbf{Fourth Fundamental Catalogue:} \label{subsubsec:fk4} The Fourth Fundamental Catalogue (FK4) defined the B1950.0 optical reference frame and constituted the backbone of mid-20th-century astrometry. Although it contains only 1,535 bright fundamental stars, FK4 provided a highly consistent realization of the reference system at the time of its construction. Many subsequent ground-based catalogs, including AGK3 and SAO, were directly or indirectly tied to FK4. As a result, its reference frame properties propagated into a large fraction of historical astrometric observations.
%     \item \textbf{Smithsonian Astrophysical Observatory Catalog:} \label{subsubsec:sao} The SAO Catalog is a compiled all-sky stellar catalog produced to supply sufficient reference star density for practical astrometric applications. It contains approximately 259,000 stars and was widely used in the early computer era for optical observation reduction, including satellite tracking. The northern portion of SAO primarily relies on AGK-based data, while the southern sky incorporates data from the Yale and Cape catalogs. This heterogeneous construction reflects the state of global astrometric resources available at the time.
%     \item \textbf{Astrographic Catalog Reference Stars:} \label{subsubsec:acrs} The Astrographic Catalog Reference Stars (ACRS) catalog was compiled by the U.S. Naval Observatory to provide reference stars for the re-reduction of the Astrographic Catalogue. It combines northern hemisphere data from AGK3 with southern hemisphere data from the Second Cape Photographic Catalogue (CPC2). ACRS was transformed onto the FK5 reference frame and represents an intermediate step between early photographic catalogs and later compiled all-sky catalogs with improved reference frame consistency.
%     \item \textbf{ACT Reference Catalog:} \label{subsubsec:act} The ACT catalog was constructed by combining early epoch positions from the Astrographic Catalogue with later observations from Tycho-1. This approach significantly improved proper motion accuracy compared to Tycho-1 alone, reaching typical values of a few milliarcseconds per year. ACT played an important role during the transition period between classical ground-based astrometry and space-based catalogs in the Hipparcos era.
%     \item \textbf{Positions and Proper Motions Catalog:} \label{subsubsec:ppm} The PPM catalog was developed as a successor to SAO, with the objective of providing a uniform realization of the J2000.0 reference frame. It contains nearly 380,000 stars and incorporates data from AGK3 in the north and FOKAT-S in the south, supplemented by early photographic observations to extend the time baseline for proper motion determination. Before the advent of CCD-based surveys, PPM represented one of the most advanced compiled ground-based astrometric catalogs.
%     \item \textbf{Yale Zone Catalogs:} \label{subsubsec:yale} The Yale Zone Catalogs consist of a series of regional photographic catalogs compiled by the Yale University Observatory. Their observations were mainly conducted between the 1930s and 1950s and focused on equatorial and southern sky regions. These catalogs were designed to improve the efficiency of stellar position measurements over wide fields and later contributed substantially to several compiled all-sky catalogs. In terms of astrometric precision, the Yale Zone Catalogs typically reach sub-arcsecond accuracy, characteristic of mid-20th-century photographic astrometry.
%     \item \textbf{Hipparcos Catalog:} \label{subsubsec:hipparcos} The Hipparcos catalog is the result of the first space astrometry mission and provides a high-precision, nearly rigid optical reference frame. It contains about 118,000 stars with typical positional and proper motion accuracies at the milliarcsecond level, with a mean observational epoch around 1991.25. Hipparcos established the foundation for all subsequent modern astrometric catalogs and serves as a fundamental link between ground-based and space-based reference systems.
%     \item \textbf{Tycho-2 Catalog:} \label{subsubsec:tycho2} The Tycho-2 catalog combines Hipparcos star mapper observations with data from 144 ground-based catalogs to produce positions and proper motions for approximately 2.5 million stars. Typical positional accuracies are of order 60 mas, with proper motion accuracies around 2–3 mas/yr. Prior to Gaia, Tycho-2 was the most widely used dense all-sky catalog for high-precision astrometric applications.
%     \item \textbf{USNO-A2.0 Catalog:} \label{subsubsec:usnoa2} USNO-A2.0 is a very large astrometric catalog generated from scans of photographic Schmidt plates, containing over 500 million sources. The catalog does not provide proper motion information; all positions correspond to the plate observation epochs, typically around 1950 for POSS-I and around 1980 for POSS-II and SRC plates. Due to its extremely high source density, USNO-A2.0 was extensively used for optical observation reduction before the availability of dense CCD-based catalogs.
%     \item \textbf{Guide Star Catalog 1.2:} \label{subsubsec:gsc12} GSC 1.2 is a re-calibrated version of the original Guide Star Catalog, developed primarily to support space telescope pointing requirements. It contains roughly 19 million stars and is based on photographic plate material, with positional accuracies at the few tenths of an arcsecond level. Although it does not include proper motions, GSC 1.2 has been widely used as an auxiliary reference catalog for ground-based astrometry.
%     \item \textbf{UCAC 2:} \label{subsubsec:ucac2} UCAC2 is the second release of the U.S. Naval Observatory CCD Astrograph Catalog and represents the transition from photographic plates to CCD-based astrometry. Its observations were conducted between approximately 1998 and 2003 and cover most of the sky from -90° to about +40° in declination. The catalog contains about 48 million stars, with positional accuracies ranging from roughly 15 to 70 mas and proper motion accuracies at the level of 1-7 mas/yr.
%     \item \textbf{UCAC 4:} \label{subsubsec:ucac4} UCAC4 is the final release of the UCAC project and provides all-sky coverage. It includes over 113 million stars and incorporates improved proper motion determinations through the use of additional early-epoch photographic material. With typical positional accuracies of about 15–20 mas for stars in the magnitude range 10–14, UCAC4 represents the mature stage of ground-based CCD astrometric catalogs prior to the Gaia era.
% \end{itemize}

\section{Basic Framework of the Correction Method} \label{sec:methods}
\subsection{Systematic Correction Model}

The systematic correction model adopted in this work follows the formalism introduced by \citet{Chesley2010} and subsequently refined by \citet{Farnocchia2015} and \citet{Eggl2020}, with Gaia DR3 serving as the astrometric reference catalog. Local stellar systematics are parameterized in terms of position offsets at epoch J2000.0, $(\Delta \alpha_{2000}, \Delta \delta_{2000})$, together with their effective linear extrapolation terms in time, $(\Delta \mu_{\alpha}, \Delta \mu_{\delta})$.

For an optical observation obtained at epoch $t$, the correction induced by reference star catalog systematics is expressed as

\begin{equation}
\Delta \alpha(t) = \Delta \alpha_{2000} + \Delta \mu_{\alpha}\,\Delta t,
\end{equation}
\begin{equation}
\Delta \delta(t) = \Delta \delta_{2000} + \Delta \mu_{\delta}\,\Delta t,
\end{equation}

where $\Delta t$ denotes the elapsed time from J2000.0 in Julian years. The correction in right ascension implicitly includes the spherical metric factor $\cos\delta$, consistent with the treatment in \citet{Chesley2010}. 

After applying these offsets, historical optical observations are effectively placed onto the Gaia DR3 reference frame and can be incorporated into orbit determination within a unified statistical framework. To ensure that the derived residuals reflect intrinsic catalog systematics rather than inconsistencies in reference definitions, all catalogs are reduced to a common ICRS/J2000 baseline prior to comparison. Catalogs already expressed in the ICRS at epoch J2000 require only epoch propagation of Gaia DR3 reference positions, whereas catalogs defined in alternative reference systems or epochs necessitate explicit frame transformations and epoch alignment before differencing.

\subsection{Catalog Harmonization: The Case of AGK1}

As an illustrative example of the harmonization procedure, the AGK1 catalog provides stellar coordinates originally defined in the FK4 reference frame at epoch B1875.0. To establish consistency with Gaia DR3, the coordinates are first transformed from FK4@B1875.0 to FK4@B1950.0 following the standard Newcomb-based FK4 conventions (e.g., \citealt{Meeus1998}), including the removal of E-terms of aberration. A subsequent reference-frame transformation is then applied to convert the coordinates from FK4 to the ICRS\citep{Murray1989}.

Gaia DR3 reference star positions are propagated backward to epoch B1950.0 using their published proper motions and are then differenced against the transformed AGK1 positions. The resulting positional and effective proper motion differences are subsequently reduced to epoch J2000.0, yielding the systematic correction parameters $(\Delta \alpha_{2000}, \Delta \delta_{2000}, \Delta \mu_{\alpha}, \Delta \mu_{\delta})$ for this catalog. 

Other historical catalogs are treated analogously, with modifications dictated by their specific reference-frame definitions and by the availability or absence of proper motion information, following the general debiasing strategy described by \citet{Farnocchia2015} and updated in the Gaia era by \citet{Eggl2020}.

\section{Results and Validation} \label{sec:results}
\subsection{Global Statistics} \label{subsec:global_corrections}

Based on the debiasing procedure described above, we derive unified estimates of the systematic position and proper motion offsets of 19 historical and modern astrometric star catalogs with respect to the reference catalog. All correction terms are propagated to epoch J2000.0 and characterized within HEALPix-defined sky regions using robust statistical estimators, namely the median (MED) and the median absolute deviation (MAD)\citep{Gorski2005}. This section first presents the global statistical properties of the derived corrections for each catalog. Representative catalogs are then examined in detail to illustrate the statistical distributions and all-sky spatial structures of the corrections, followed by an analysis of their systematic origins and a validation of the effectiveness of the applied smoothing procedures.

\begin{deluxetable*}{ccccccccc}
\tablecaption{Global statistics of position and proper motion corrections for each star catalog at epoch J2000.0\label{tab:global_corrections_j2000}}
\tablehead{
\colhead{Catalog} & \multicolumn{2}{c}{$\Delta\alpha_{J2000}$ [mas]} & \multicolumn{2}{c}{$\Delta\delta_{J2000}$ [mas]} & \multicolumn{2}{c}{$\Delta\mu_{\alpha}$ [mas/yr]} & \multicolumn{2}{c}{$\Delta\mu_{\delta}$ [mas/yr]} \\
\colhead{} & \colhead{MED} & \colhead{MAD} & \colhead{MED} & \colhead{MAD} & \colhead{MED} & \colhead{MAD} & \colhead{MED} & \colhead{MAD}}
\startdata
AC            & $-128.2$ & $257.6$ & $-44.9$  & $301.4$ & $-1.82$ & $4.67$ & $-0.58$ & $5.37$ \\
ACRS          & $-39.3$  & $136.9$ & $-48.5$  & $128.1$ & $-0.90$ & $2.86$ & $-0.07$ & $2.56$ \\
ACT           & $-0.4$   & $11.6$  & $0.2$    & $10.8$  & $-0.08$ & $1.03$ & $0.01$  & $0.97$ \\
AGK1          & $-128.2$ & $257.6$ & $-44.9$  & $301.4$ & $-1.82$ & $4.67$ & $-0.58$ & $5.37$ \\
AGK3          & $-128.2$ & $257.6$ & $-44.9$  & $301.4$ & $-1.82$ & $4.67$ & $-0.58$ & $5.37$ \\
FK4 Catalogue & $-73.1$  & $157.7$ & $-46.7$  & $200.7$ & $-0.89$ & $2.32$ & $-0.28$ & $3.03$ \\
GSC 1.2       & $28.2$   & $113.2$ & $-13.5$  & $104.4$ & $1.03$  & $2.89$ & $4.00$  & $1.52$ \\
Gaia DR1      & $-23.7$  & $34.7$  & $-56.2$  & $23.3$  & $1.57$  & $2.31$ & $3.74$  & $1.55$ \\
Gaia DR2      & $-0.0$   & $0.3$   & $-0.2$   & $0.4$   & $0.00$  & $0.02$ & $0.01$  & $0.02$ \\
Hipparcos     & $0.7$    & $6.6$   & $0.1$    & $5.6$   & $0.09$  & $0.72$ & $0.01$  & $0.61$ \\
PPM           & $12.7$   & $128.4$ & $-67.7$  & $127.1$ & $-0.23$ & $2.71$ & $-0.45$ & $2.77$ \\
SAO           & $-5.8$   & $270.6$ & $-19.4$  & $266.9$ & $0.62$  & $3.73$ & $-0.07$ & $3.63$ \\
Tycho-2       & $-0.4$   & $10.0$  & $-1.1$   & $9.7$   & $-0.07$ & $0.81$ & $-0.25$ & $0.81$ \\
UCAC2         & $-0.8$   & $10.3$  & $2.8$    & $9.2$   & $-0.29$ & $1.31$ & $-0.57$ & $1.25$ \\
UCAC4         & $1.0$    & $11.0$  & $-3.0$   & $9.6$   & $0.11$  & $0.84$ & $0.37$  & $0.94$ \\
USNO-A2.0     & $68.6$   & $107.6$ & $156.8$  & $124.6$ & $1.36$  & $1.92$ & $3.31$  & $1.32$ \\
Yale          & $19.3$   & $410.1$ & $-139.9$ & $439.7$ & $0.58$  & $5.70$ & $-0.63$ & $5.74$ \\
\enddata
\tablecomments{The tabulated values correspond to the median (MED) and median absolute deviation (MAD) of the correction terms over all valid HEALPix sky regions. Position corrections refer to the right ascension and declination components at epoch J2000.0, where $\Delta\alpha$ includes the $\cos\delta$ factor. Proper motion corrections are given in units of mas\,yr$^{-1}$.}
\end{deluxetable*}

Table~\ref{tab:global_corrections_j2000} summarizes the global robust statistics of position and proper motion corrections for each star catalog at epoch J2000.0. Overall, the amplitudes of systematic offsets exhibit a clear stratification among different catalogs. Modern CCD-based catalogs (e.g., Gaia DR2, Hipparcos, and Tycho-2) show global medians and dispersions at the milliarcsecond level, whereas historical photographic catalogs and catalogs lacking proper motion information typically display position systematics exceeding $10^2$\,mas, accompanied by substantial spatial scatter.

Taking USNO-A2.0 as an illustrative example, the global medians of $\Delta\alpha_{\mathrm{J2000}}$ and $\Delta\delta_{\mathrm{J2000}}$ reach $68.6$\,mas and $156.8$\,mas, respectively, with corresponding MAD values exceeding $100$\,mas. These results indicate that such catalogs deviate significantly from the accuracy requirements of modern astrometric reference frames when evaluated in the J2000 context.

\subsection{Statistical distributions and spatial structures of the corrections} \label{subsec:distribution}

Figure~\ref{fig:hist_usnoa2} presents the distributions of the accumulated correction terms for the USNO-A2.0 catalog over all HEALPix sky regions. The position corrections exhibit pronounced global offsets in both right ascension and declination, with distributions that deviate significantly from a Gaussian shape. The proper motion corrections also display systematic biases, with the declination component showing a larger median amplitude than the right ascension component. This behavior reflects an asymmetric amplification effect introduced during epoch propagation as a consequence of the absence of intrinsic proper motion information in the catalog.

\begin{figure}[ht!]
\plotone{../Figue/fig_bias_usno_a2_j2000.png}
% \caption{USNO-A2.0 Summary statistics of USNO A2.0 catalog systematics with respect to Gaia DR 3}
\caption{Statistical distributions of the correction terms for the USNO-A2.0 star catalog}
\label{fig:hist_usnoa2}
\end{figure}

Beyond their statistical distributions, the correction terms exhibit pronounced spatially correlated structures across the sky. As shown in Figure~\ref{fig:bias_usno_a2_2x2}, the position and proper motion corrections for the USNO-A2.0 catalog display smoothly varying gradient-like patterns over multiple sky-region scales, together with distinct block-like discontinuities in localized areas. The presence of such spatial correlations indicates that the derived corrections predominantly trace intrinsic systematic error patterns of the star catalog, rather than being driven by random statistical fluctuations associated with cross-matching noise.

\begin{figure}[ht!]
\plotone{../Figue/bias_usno_a2_2x2.png}
\caption{All-sky spatial distributions of the correction terms for the USNO-A2.0 star catalog}
\label{fig:bias_usno_a2_2x2}
\end{figure}

\subsection{Interpretative analysis of systematic bias characteristics} \label{subsec:interpretation}

The spatial structures exhibited by the correction terms of different star catalogs can be interpreted within a unified framework, as they are primarily governed by the original observational techniques, reference frame realizations, and epoch propagation strategies. Consistent with the findings of \citet{Eggl2020}, we find that the systematic offsets of historical star catalogs relative to the Gaia reference frame are not randomly distributed, but are instead closely linked to their respective observation and reduction methodologies.

Photographic catalogs such as GSC~1.2, AGK3, SAO, and the Yale catalog exhibit pronounced block-like spatial structures across the sky. These spatial discontinuities directly correspond to the coverage of individual photographic plate fields and their independent reduction procedures\citep{Platais1998,Eichhorn1974}. Further quantitative analysis shows that significant jumps in the correction terms are ubiquitous at the boundaries between adjacent HEALPix regions for this class of catalogs. For the Yale, AGK3, and SAO catalogs, more than 97\% of sky-region boundaries show correction jumps $\delta_{\text{jump}}(i,j)$ exceeding $100$\,mas. In the case of the Yale catalog, approximately 45\% of the boundaries exhibit $\delta_{\text{jump}}(i,j)$ larger than $1''$, with maximum values reaching several arcseconds (Yale $\sim7.3''$, AGK3 $\sim3.8''$). These results indicate that such discontinuities do not arise from a small number of anomalous sky regions, but rather represent a global manifestation of systematic distortions and fitting residuals inherited from independent plate-field reductions and regional stitching procedures. This behavior highlights the intrinsic limitations of photographic catalogs that lack a global block adjustment based on all-sky overlapping constraints.

\begin{equation}
\delta_{\mathrm{jump}}(i,j) = \left| V(i) - V(j) \right|
\end{equation}

where:

\begin{itemize}
\item $i$ and $j$ denote two spatially adjacent HEALPix pixels on the celestial sphere, i.e., $j \in \mathrm{Neighbors}(i)$;
\item $V(\cdot)$ represents the value of the correction term evaluated at the corresponding pixel (e.g., $\Delta\alpha^*$ or $\Delta\delta$);
\item $\delta_{\mathrm{jump}}$ quantifies the amplitude of the discontinuity between the two neighboring pixels.
\end{itemize}

For star catalogs based on scanning or drift-scan observing strategies (e.g., UCAC2), systematic biases typically manifest as stripe-like structures aligned with the scan direction\citep{Zacharias2017UCAC5}, with correction amplitudes increasing significantly in regions of low Galactic latitude and high source density. This behavior is plausibly associated with charge transfer efficiency effects in early CCD detectors and with measurement degradation in crowded stellar fields. By contrast, UCAC4 incorporates higher-quality reference catalogs and an improved reduction pipeline, which substantially mitigates systematics correlated with Galactic structure.

For catalogs lacking proper motion information (e.g., USNO-A2.0), the dominant systematic signature appears as a smooth, large-scale gradient across the sky. Small systematic errors in stellar proper motions are linearly amplified over several decades of epoch propagation\citep{Vasiliev2019}, producing pronounced artificial patterns in a static coordinate frame. Similar behavior can also be identified in Gaia DR1, further underscoring the critical role of accurate proper motion measurements in establishing a rigid astrometric reference frame.

\subsection{Smoothing of the correction results} \label{subsec:smoothing}

Because several star catalogs exhibit pronounced discontinuities in the correction fields at HEALPix tile boundaries, directly applying discrete tile-based corrections may introduce non-physical coordinate jumps in the vicinity of these boundaries, thereby degrading the numerical stability of orbit determination and residual analysis. Consequently, an engineering-level spatial smoothing of the correction fields is required, provided that the large-scale spatial structures of the corrections are preserved.

\begin{figure}[ht!]
\plotone{"../Figue/bias_ucac4_zoom_compare_n256.png"}
\caption{Basis Function Smoothing}
\label{fig:bias_ucac4_zoom_compare_n256}
\end{figure}

In this work, we adopt the Basis Function Smoothing (BFS) approach proposed by \citet{Eggl2020} to smooth the discrete correction fields defined on the HEALPix grid\citep{Wahba1990,Silverman1986}. This procedure operates solely on the spatial continuity of the correction fields and does not alter their global statistical properties. The smoothed corrections are evaluated and represented on a higher-resolution grid with \texttt{nside}=256, which effectively suppresses boundary-induced discontinuities while preserving the large-scale systematic error patterns.

\begin{figure*}[ht!]
\plotone{"../Figue/bias_gsc1.2_smooth_compare_allsky_2x3_v2.png"}
\caption{BFS-smoothed correction fields at \texttt{nside}=256}
\label{fig:bias_smooth}
\end{figure*}

The impact of smoothing on field continuity is quantitatively assessed using the inter-pixel discontinuity metric between neighboring HEALPix regions. The results show that BFS smoothing significantly reduces the amplitudes of boundary discontinuities for all analyzed catalogs, with both the 95th percentile and the maximum values decreasing to approximately $10\%$ of those in the original correction fields.

\begin{figure}[ht!]
\plotone{"../Figue/figure_Y_jump_distribution.png"}
\caption{Cumulative distributions of inter-pixel discontinuities $\delta_{\mathrm{jump}}(i,j)$ before and after smoothing}
\label{fig:bias_jump}
\end{figure}

At the same time, the global statistical properties of the correction fields remain effectively unchanged after smoothing. Variations in the MED and MAD are substantially smaller than the intrinsic dispersion of the corresponding distributions\citep{Cleveland1979}. This demonstrates that BFS smoothing does not significantly modify the overall amplitudes or the large-scale spatial structures of the corrections, but instead primarily mitigates numerical discontinuities at tile boundaries, thereby improving the continuity and numerical robustness of the correction model in practical applications.

\subsection{Validation} \label{subsec:validation}
\subsection{Plate Re-reduction Validation}

To independently assess the effectiveness of the catalog systematic correction model, we adopt a plate re-reduction benchmark based on two independent reductions of the same Triton CCD frames obtained during 1996--2006. Qiao et al. (2007) reduced these observations using the UCAC2 star catalog, providing 943 astrometric positions \citep{Qiao2007}. Yan et al. (2020) subsequently re-measured and re-reduced the identical CCD images within the Gaia DR2 reference frame, yielding 1007 positions including 941 updated measurements \citep{Yan2020}. Because both solutions share the same original CCD material, the dominant difference between them arises from the adopted reference star catalog and astrometric frame realization.

From the overlapping data, 930 strictly corresponding observations were selected. The Gaia DR2 re-reduction solution is adopted as the high-precision reference. The total positional separation is defined as

\begin{equation}
\Delta = \sqrt{(\Delta\alpha \cos\delta)^2 + (\Delta\delta)^2},
\end{equation}

where the right ascension component explicitly includes the spherical metric factor $\cos\delta$.

Prior to correction, the UCAC2-based solution exhibits a mean separation of 86.2 mas (median 82.5 mas) relative to the Gaia DR2 re-reduction. After applying the Gaia DR3-based regional systematic correction model developed in this work, the mean separation decreases to 76.2 mas (median 67.7 mas), corresponding to an overall improvement of 10.5\%. A total of 78.0\% of individual observations show reduced residuals.

\begin{figure}[ht!]
\plotone{../Figue/fig_plate_validation.png}
\caption{(a) Cumulative distributions of total positional separation $\Delta$ before (red) and after (blue) correction. (b) Temporal evolution of residuals: moving-average trends (upper) and net gain $\Delta_{\mathrm{pre}}-\Delta_{\mathrm{post}}$ (lower).}
\label{fig:plate_validation}
\end{figure}

Figure~\ref{fig:plate_validation}a presents the cumulative distribution functions (CDFs) of $\Delta$ before and after correction. The post-correction curve is globally shifted toward smaller separations across the full range, indicating that the improvement is not confined to the central tendency but affects the entire residual distribution\citep{Veres2017}. 

Figure~\ref{fig:plate_validation}b shows the temporal stability analysis. The upper panel displays moving-average trends of residuals before and after correction, while the lower panel presents the point-wise net gain ($\Delta_{\mathrm{pre}} - \Delta_{\mathrm{post}}$). The reduction persists throughout the full observational interval and is not restricted to specific epochs, demonstrating that the improvement is statistically consistent and temporally stable.

Although the correction model is constructed using Gaia DR3 while the validation reference is based on Gaia DR2, no methodological inconsistency is introduced. Both releases share the same ICRS realization, and DR3 provides an extended observational time baseline and improved astrometric parameter solutions \citep{GaiaDR2,GaiaDR3}. The internal differences between DR3 and DR2 are substantially smaller than the systematic offset between UCAC2 and the Gaia reference frame. Therefore, convergence toward the Gaia DR2 solution constitutes a conservative and stringent validation of the DR3-based correction model.

\subsection{Ephemeris Comparison Validation}

To assess the impact of the systematic error correction model at the dynamical level, we perform an independent external validation using ephemeris comparison. We select Jovian satellite astrometric observations published by IMCCE, covering samples from various historical catalogs including ACT, AGK3, and UCAC4. All observation times are converted to the Terrestrial Time (TT) scale, and coordinates are unified to the ICRS reference frame. Theoretical positions obtained from JPL Horizons serve as the dynamical reference. Residuals with respect to the ephemeris are computed for each observation both before and after applying the corrections. The statistical results are summarized in Table~\ref{tab:ephem_validation}, where R1 denotes the statistics before correction and R2 denotes those after correction.

Regarding the mean residuals, several samples show a convergence toward zero in Right Ascension after correction. For instance, in the ACT dataset, the mean RA residual for sample jg0006 improves from $-0.0297''$ to $-0.0181''$, and for jo0001 from $0.0058''$ to $-0.0023''$; similarly, for the UCAC4 sample jg0045, the mean reduces from $0.0205''$ to $0.0118''$. These variations suggest that the correction model effectively mitigates local systematic offsets.

In terms of standard deviation, the width of the residual distribution remains largely stable across samples. Taking the large UCAC4 sample jo0056 ($N=3613$) as an example, the RA standard deviation shifts marginally from $0.6437''$ to $0.6447''$, and in Declination from $9.1106''$ to $9.1117''$. Similar behavior is observed for the ACT and AGK3 samples, indicating that the correction process does not significantly alter the dispersion of the residuals.

Figures~\ref{fig:vector_ucac4}--\ref{fig:vector_agk3} display the planar distribution of the residual vectors. The overall scatter ranges remain consistent before and after correction, with no emergence of new systematic banding or directional biases, confirming that the correction procedure introduces no additional dynamical systematic errors.

In summary, the ephemeris comparison validates that the Gaia DR3-based systematic error model constructed in this study preserves the stochastic error structure of the observational data at the dynamical level while partially reducing systematic mean offsets in specific samples.

\begin{deluxetable*}{cccccccccc}
\tablecaption{Ephemeris residual statistics before (R1) and after (R2) catalog bias correction. Units are arcseconds.
\label{tab:ephem_validation}}
\tablehead{\colhead{Catalog} & \colhead{N} &\colhead{R1 RA Mean} & \colhead{R2 RA Mean} &\colhead{R1 RA Std} & \colhead{R2 RA Std} &\colhead{R1 Dec Mean} & \colhead{R2 Dec Mean} &\colhead{R1 Dec Std} & \colhead{R2 Dec Std}}
\startdata
ACT (\href{https://www.sai.msu.ru/neb/nss/html/obspos/OBS_COLL/J/jg0005.html}{jg0005}) & 218 & -0.0084 & -0.0126 & 0.1158 & 0.1150 & -0.0316 & -0.0316 & 0.1567 & 0.1548 \\
ACT (\href{https://www.sai.msu.ru/neb/nss/html/obspos/OBS_COLL/J/jg0006.html}{jg0006}) & 224 & -0.0297 & -0.0181 & 0.0861 & 0.0876 & -0.0451 & -0.0478 & 0.1474 & 0.1449 \\
ACT (\href{https://www.sai.msu.ru/neb/nss/html/obspos/OBS_COLL/J/jo0001.html}{jo0001}) & 115 & 0.0058 & -0.0023 & 0.1907 & 0.1899 & -0.0163 & -0.0118 & 0.1872 & 0.1892 \\
AGK3 (\href{https://www.sai.msu.ru/neb/nss/html/obspos/OBS_COLL/J/ji0007.html}{ji0007}) & 10 & 0.2939 & 0.2708 & 0.1715 & 0.1715 & 0.0112 & -0.0676 & 0.0704 & 0.0704 \\
UCAC4 (\href{https://www.sai.msu.ru/neb/nss/html/obspos/OBS_COLL/J/jg0045.html}{jg0045}) & 154 & 0.0205 & 0.0118 & 0.0773 & 0.0769 & -0.0193 & -0.0188 & 0.0848 & 0.0864 \\
UCAC4 (\href{https://www.sai.msu.ru/neb/nss/html/obspos/OBS_COLL/J/jo0056.html}{jo0056}) & 3613 & 0.5442 & 0.5481 & 0.6437 & 0.6447 & 1.9799 & 1.9805 & 9.1106 & 9.1117 \\
\enddata
\end{deluxetable*}

\begin{figure}[ht!]
\plotone{../Figue/UCAC4_best_vector.png}
\caption{Residual vector distribution for UCAC4 observations.}
\label{fig:vector_ucac4}
\end{figure}
\begin{figure}[ht!]
\plotone{../Figue/ACT_best_vector.png}
\caption{Residual vector distribution for ACT observations.}
\label{fig:vector_act}
\end{figure}
\begin{figure}[ht!]
\plotone{../Figue/AGK3_best_vector.png}
\caption{Residual vector distribution for AGK3 observations.}
\label{fig:vector_agk3}
\end{figure}

\section{Conclusions} \label{sec:conclusions}

Based on Gaia DR3, which currently provides the highest-precision astrometric reference frame, we present systematic position and proper motion correction tables for 17 historical astrometric star catalogs. A unified debiasing framework is constructed by combining HEALPix-based sky partitioning with Basis Function Smoothing (BFS), enabling the characterization and suppression of spatial discontinuities in catalog systematics across the entire sky. The results demonstrate that this approach effectively smooths high-frequency numerical artifacts introduced by discrete sky partitioning while preserving the genuine large-scale spatial structures of catalog systematics.

Through comparative validation of astrometric observations before and after applying the proposed corrections, we find that the corrected data exhibit reduced residuals in a statistical sense, together with improved consistency and stability in orbit determination solutions\citep{Desmars2015}. This indicates that the correction model efficiently mitigates the impact of reference catalog systematics on orbit determination, and is particularly beneficial for historical observation epochs that lack direct Gaia-based reference stars.

The correction tables provided in this work are intended as an a posteriori debiasing tool for the processing of historical astrometric data, aimed at improving the statistical reliability of long-arc orbit determinations. For contemporary and future astrometric observations, however, optimal accuracy is best achieved by directly adopting the Gaia star catalog and its subsequent data releases as the astrometric reference.

\section{Data and Code Availability} \label{sec:data}

The systematic bias correction tables derived in this work are provided at two resolutions to accommodate different precision requirements: a standard-resolution version (NSIDE=64) and a smoothed high-resolution version (NSIDE=256). Detailed documentation is included with each dataset.

The correction tables are publicly available at \url{https://github.com/yangzhe0/nside64} and \url{https://github.com/yangzhe0/nside256}.

In addition, the source code for the bias estimation pipeline used in this work is openly released to facilitate reproducibility and further analysis. The code repository can be accessed at \url{https://github.com/yangzhe0/catalog}.

\begin{acknowledgments}
    VizieR catalogue access tools, CDS, Strasbourg, France, have been essential to this research project (Ochsenbein et al., 2000). This work made use of the IAU Minor Planet Center observations database.
\end{acknowledgments}

\begin{contribution}
%%This section gives authors the space to recognize author contributions. The text inside this environment is NOT counted towards the total word quanta. At a minimum, manuscripts are expected to include this text:
All authors contributed equally to the Terra Mater collaboration.
%% But authors are expected to provide more specific details, e.g. 
%%
%%SC was responsible for writing and submitting the manuscript.
%%WWM came up with the initial research concept and edited the manuscript.
%%OTS obtained the funding and edited the manuscript.
%%EBF provided the formal analysis and validation. He also edited the manuscript.
%%GEH Supervised the undergraduates, wrote the software and administers the project github and Zenodo repositories.
%%
%% Authors can use the Contributor Role Taxonomy (CRediT) at
%% https://credit.niso.org
%% for ideas on how write a good statement tailored to their needs.
\end{contribution}

%% Similar to \facility{}, there is the optional \software command to allow 
%% authors a place to specify which programs were used during the creation of 
%% the manuscript. Authors should list each code and include either a
%% citation or url to the code inside ()s when available.
\software{HEALPix \citep{Gorski2005}, Astropy \citep{Astropy2013}} 

\bibliography{reference}{}
\bibliographystyle{aasjournalv7}
\end{document}
